\documentclass[UTF8,11pt,a4paper,fontset=none]{ctexart}

\usepackage[a4paper,margin=2.25cm]{geometry}
\usepackage{fontspec}
\usepackage{xeCJK}
\setmainfont{DejaVu Serif}
\setsansfont{DejaVu Sans}
\setmonofont{DejaVu Sans Mono}
\setCJKmainfont{Noto Serif CJK SC}
\setCJKsansfont{Noto Sans CJK SC}
\setCJKmonofont{Noto Sans Mono CJK SC}
\usepackage{amsmath,amssymb}
\usepackage{booktabs,longtable,tabularx,array,multirow}
\usepackage[table]{xcolor}
\usepackage{enumitem}
\usepackage{hyperref}
\usepackage{xurl}
\usepackage{microtype}
\usepackage{fancyhdr}
\usepackage{caption}

\definecolor{NimlothBlue}{HTML}{28557A}
\definecolor{NimlothLight}{HTML}{EEF4F8}
\definecolor{NimlothGreen}{HTML}{2E6B4F}
\definecolor{NimlothRed}{HTML}{A33A3A}
\hypersetup{
  colorlinks=true,
  linkcolor=NimlothBlue,
  urlcolor=NimlothBlue,
  citecolor=NimlothBlue,
  pdftitle={Nimloth近期状态接口与世界模型调查报告}
}
\captionsetup{font=small,labelfont=bf}
\setlist{leftmargin=2em,itemsep=0.25em,topsep=0.35em}
\renewcommand{\arraystretch}{1.18}
\setlength{\parindent}{2em}
\setlength{\parskip}{0.25em}
\pagestyle{fancy}
\fancyhf{}
\fancyhead[L]{Nimloth近期调查报告}
\fancyhead[R]{2026年8月}
\fancyfoot[C]{\thepage}

\newcommand{\good}[1]{\textcolor{NimlothGreen}{\textbf{#1}}}
\newcommand{\bad}[1]{\textcolor{NimlothRed}{\textbf{#1}}}
\newcommand{\code}[1]{\texttt{#1}}
\newcommand{\R}{\mathbb{R}}

\title{\textbf{Nimloth近期状态接口与世界模型调查报告}\\[0.5em]
\large 从长期规划失败到视觉与任务信息的定位}
\author{Nimloth项目}
\date{2026年8月24日}

\begin{document}
\maketitle

\begin{abstract}
Nimloth希望让智能体先在内部预测动作后果，再选择真实动作。近期评估表明，当前系统已经能够完成一部分短任务，但在长期任务上明显不足。调查随后集中到一个基础问题：模型交给世界模型的内部状态，是否真的保留了规划所需的画面、任务目标和动作可行性信息？

本报告汇总ID54--ID75、ID191和ID192等近期调查。主要结论有三点。第一，部署中的语言与视觉主模型本身基本健康，严重漂移主要发生在旧训练阶段的状态压缩模块。第二，当前16位置状态几乎没有保留任务目标，并削弱了部分左右碰撞信息；世界模型又进一步没有充分利用剩余信息，尤其容易把被障碍挡住的移动预测成成功移动。第三，最新调查已经找到更可靠的信息来源：任务目标在原始任务指令（instruction）表示中几乎可被完全读出；经过语言模型处理后的当前图片表示，在三个主要移动方向上接近DINO视觉基准，但左右动作的外部样本仍不足以通过严格统计门槛。

因此，下一步不应继续修补旧世界模型或扩大旧适配器。本报告进一步对比了当前SFT1监督与建议的新监督：当前训练只有成功示范的回答/动作损失和把状态直接拟合到DINO的视觉损失，缺少任务指令、动作被挡和目标--画面关系监督。更稳妥的路线是重新监督SFT1的16位置状态提取；只有新状态同时通过视觉、目标、动作可行性和策略保持检查后，才从头训练新的世界模型、价值模型和搜索策略。
\end{abstract}

\tableofcontents
\newpage

\section{报告范围与核心问题}

本轮调查由正式规划评估中的两个现象触发：

\begin{itemize}
  \item 训练到第20次更新后的K4搜索策略，在300个测试任务上成功率为$110/300=36.67\%$；
  \item 分类型看，长期任务只有$2/60$成功，远低于其他类型。
\end{itemize}

随后完成的Base与Common Sense各60条只读评估中，成功数分别为$26/60$和$29/60$。这些结果说明系统并非完全失效，但搜索深度增加后并没有稳定带来长期能力。为了判断原因，我们依次调查了以下问题：

\begin{enumerate}
  \item 世界模型预测的内部状态，是否仍处在真实状态通常出现的范围内？
  \item 失败来自部署中的主模型，还是来自把主模型输出压缩成16位置状态的模块？
  \item 当前状态是否保留了任务目标？
  \item 当前状态是否能判断向前、向左或向右移动会成功还是被挡住？
  \item 如果信息丢失，它是在视觉编码、语言与图像融合、instruction表示，还是最后的16位置压缩阶段丢失？
\end{enumerate}

本报告重点回答这些问题。ValueHead在第2--4层预测状态上的完整校准调查尚未完成，因此不在本报告中给出确定结论。

\section{系统中“状态”的含义}

每个真实时刻，智能体会看到当前图片、任务instruction和此前交互，并生成当前回答。项目使用这次真实生成得到的内部表示，压缩成16个位置，每个位置有1024个数：

\begin{equation}
  z_t \in \R^{16\times 1024}.
\end{equation}

报告中将$z_t$称为“16位置状态”。它应当是对当前局面的简洁描述，至少保留：

\begin{itemize}
  \item 当前画面中的空间与障碍信息；
  \item instruction指定的任务目标；
  \item 目标与当前画面的关系；
  \item 对动作后果有用的信息。
\end{itemize}

世界模型根据当前状态和动作预测下一状态：

\begin{equation}
  \hat z_{t+1}=F(z_t,a_t).
\end{equation}

价值模型再读取真实或预测状态，估计动作的长期结果。搜索会重复使用世界模型和价值模型，因此第一层很小的状态偏差可能在第2--4层被放大。

DINO在调查中只作为视觉比较基准。它擅长描述图片内容，但不知道instruction，因此不能独自定义完整状态。

\section{调查方法与解释边界}

\subsection{数据和冻结范围}

主要状态调查使用强化学习之前的归档数据：3,211条训练轨迹和355条外部验证轨迹，并重点使用每条轨迹前4步。所有需要回答内容的状态都使用该观察对应的真实归档回答，没有填入固定、示例或人为编写的思考文本。

除明确说明的诊断模型外，部署主模型、视觉模块、状态压缩模块、DINO、旧世界模型、价值模型和强化学习模块均保持固定。训练的小型分类器只用于回答“某种信息能否从特征中读出”，不能作为部署模型或后续checkpoint使用。

\subsection{数据身份限制}

归档中的部分\code{config\_id}、seed和UID不能可靠对应到每一条实际轨迹。因此，目标标签只从归档中的完整instruction与全局一致的目标类别映射得到。我们按原始图片内容分组，排除训练与验证间完全相同的图片，但不能把结果描述为正式的逐任务泛化评估。

\subsection{主要指标如何理解}

\paragraph{RMSE。}
表示预测与目标之间的平均误差，单位与被比较的内部表示一致；越小越好。它适合比较同一类表示，但不同表示若尺度不同，RMSE不能单独决定哪一种语义更好。

\paragraph{AUC。}
用于判断一次移动会成功还是被挡住。$0.5$相当于随机猜测，$1.0$表示完全区分。AUC只衡量排序能力，不代表概率已经校准准确。

\paragraph{目标分类准确率。}
尝试从一种内部表示判断instruction中的目标类别。micro准确率按样本统计，macro准确率先分别计算每类准确率再平均，因此更不容易被常见类别主导。

\paragraph{相对“保持不变”基准的改善。}
世界模型的一步预测与“直接认为下一状态等于当前状态”比较：

\begin{equation}
  \mathrm{skill}=1-
  \frac{\operatorname{MSE}(\hat z_{t+1},z_{t+1})}
       {\operatorname{MSE}(z_t,z_{t+1})}.
\end{equation}

下文把这个数简称为skill（相对复制改善）。$\mathrm{skill}>0$表示世界模型优于“什么都没变”，小于0表示还不如直接复制当前状态。

\section{调查过程概览}

\begin{longtable}{p{1.7cm}p{4.7cm}p{8.0cm}}
\toprule
\textbf{调查} & \textbf{问题} & \textbf{主要发现} \\
\midrule
\endfirsthead
\toprule
\textbf{调查} & \textbf{问题} & \textbf{主要发现} \\
\midrule
\endhead
ID56--57 & 世界模型预测是否与真实状态、原图视觉一致？ & 预测状态更接近DINO视觉表示，却严重偏离真实16位置状态；真实状态和预测状态没有处在同一稳定尺度。 \\
ID58--59 & 漂移来自主模型还是状态压缩模块？ & 主模型变化很小；更换旧状态压缩模块会造成约$0.89$的状态RMSE变化，是主模型变化的40倍以上。 \\
ID60 & 现有状态是否保留任务目标？ & 目标分类准确率低于DINO和多数类基准，目标门槛全部失败。 \\
ID75 & 在固定状态上重新训练一步世界模型是否可行？ & 总体改善明显，但向左动作低于“保持不变”基准，因此不能继续训练更长预测。 \\
ID61 & 为什么向左预测失败？ & 成功移动可以预测；被挡住时，模型会错误预测画面发生移动。验证集中被挡比例又更高。 \\
ID71 & 碰撞信息在状态中是否存在？ & 状态保留一部分，但左右方向弱于DINO；世界模型又低于直接读取状态的小分类器。 \\
ID191 & 状态压缩之前的最后隐藏表示能否用小适配器修复？ & 不能。适配器略改善视觉距离，但没有恢复目标和左右碰撞信息。 \\
ID192 & 目标和视觉信息究竟在哪里？ & instruction表示几乎完整保留目标；语言模型处理后的当前图片表示接近DINO，是最有希望的视觉来源。 \\
\bottomrule
\end{longtable}

\section{真实状态与世界模型预测不在同一坐标中}

\subsection{ID56：预测状态远离真实状态}

在1,742个有精确下一真实状态的转移中：

\begin{itemize}
  \item 世界模型预测到真实下一状态的RMSE均值为$0.526$；
  \item 直接复制当前状态的RMSE只有$0.119$；
  \item 1,742个样本中，没有一个预测在该指标上优于复制；
  \item 真实状态的标准差约为$0.49$--$0.52$，预测状态约为$0.824$--$0.826$。
\end{itemize}

这些数字说明预测并非只是“稍微不准”。预测状态的尺度已经明显离开真实状态常见的范围。旧图像解码器即使输入真实状态也几乎不能恢复图片，因此解码图不能作为判断状态质量的可靠依据。

\subsection{ID57：预测更像DINO，但不像真实状态}

使用原始观察图片重新计算DINO后，结论仍然成立：

\begin{itemize}
  \item 预测状态到下一图片DINO的RMSE为$0.831$，优于复制的$0.981$；
  \item 但预测到真实下一状态的相对复制能力为$-9.45$；
  \item 固定地交换16个位置只能改善$0.169\%$，因此主要问题不是位置排列错误。
\end{itemize}

最合理的解释是：旧世界模型受到DINO损失影响，被拉向一种较强的视觉表示；真实状态压缩模块却输出另一种尺度和方向。价值模型主要在真实状态上学习，却在搜索深处读取预测状态，这为多层搜索带来明显风险。

\section{主模型基本健康，旧状态压缩模块发生严重漂移}

ID58比较了SFT1、旧SFT2和不同压缩模块的组合。

\begin{table}[htbp]
\centering
\caption{主模型变化与状态压缩模块变化}
\small
\begin{tabularx}{\textwidth}{Xrr}
\toprule
\textbf{比较} & \textbf{状态RMSE变化} & \textbf{解释} \\
\midrule
固定SFT1压缩模块，只把SFT1主模型换成ID74主模型 & $0.0208$ & 主模型变化很小 \\
固定ID176主模型，把SFT1压缩模块换成ID74压缩模块 & $0.8853$ & 压缩模块造成巨大变化 \\
ID58各组合中的压缩模块更换 & $0.8920$--$0.8927$ & 约为主模型变化的40倍以上 \\
\bottomrule
\end{tabularx}
\end{table}

SFT1压缩模块的视觉表现也明显更好：它到DINO的RMSE约$0.838$、余弦相似度约$0.656$；ID74压缩模块对应约$1.137$和$0.382$。ID59在实际部署主模型ID176上复现了同样结果，并确认ID176与SFT1主模型的状态差异只有$0.0188$。

因此，近期调查没有发现部署主模型整体损坏的证据。旧SFT2训练后的状态压缩模块是旧接口失效的主要来源。不过，这不表示SFT1状态已经足够好；后续目标和动作结果调查证明它仍缺少规划需要的信息。

\section{现有16位置状态几乎没有保留任务目标}

ID60在冻结状态上训练了一个容量很小的目标分类器，并用相同方法比较DINO。

\begin{table}[htbp]
\centering
\caption{ID60目标分类结果}
\small
\begin{tabular}{lcc}
\toprule
\textbf{输入表示} & \textbf{micro top-1} & \textbf{macro top-1} \\
\midrule
现有16位置状态 & $5.78\%$ & $4.04\%$ \\
DINO视觉表示 & $10.69\%$ & $6.71\%$ \\
多数类基准 & $7.23\%$ & --- \\
\bottomrule
\end{tabular}
\end{table}

状态结果低于DINO，也低于始终猜测最常见目标的基准。状态减去DINO准确率的95\%区间为$[-8.09,-1.73]$个百分点。由此不能接受现有状态作为完整的“画面加任务目标”接口。

DINO只看图片，因此它的目标分类能力可能来自物体共现或数据偏差。它高于当前状态并不表示DINO已经理解instruction；它只说明当前状态没有展现出更强、更可靠的目标信息。

\section{一步世界模型能学习移动，但无法可靠处理被挡住的动作}

\subsection{ID75：总体进步掩盖了向左失败}

我们固定状态接口，重新训练了一个从“直接复制当前状态”开始的一步世界模型。外部验证共1,413个转移：

\begin{itemize}
  \item 总体skill为$+0.365$；
  \item 主要动作平均skill为$+0.210$；
  \item 预测与真实状态标准差之比为$0.981$；
  \item 但是向前、向右平移、向左平移、向右转向动作的skill分别为
  $+0.354$、$+0.054$、$-0.133$、$+0.566$。
\end{itemize}

向左动作没有超过复制基准，所以按照预先规定的逐动作门槛，调查在一步阶段停止，没有继续训练第2--4步世界模型。

\subsection{ID61：失败集中在“动作被挡住”}

归档中的下一条环境反馈可以直接说明每次移动是否成功。向左移动在训练集和外部验证集中的失败率分别为：

\begin{equation}
  22.68\% \quad\text{和}\quad 38.86\%.
\end{equation}

外部验证多了$16.18$个百分点的失败动作。分开计算后：

\begin{itemize}
  \item 成功向左时，世界模型skill为$+0.163$；
  \item 被挡住时，skill为$-9.914$；
  \item 被挡住的真实图片$100\%$保持不变；
  \item 世界模型预测的状态变化只能以AUC $0.617$区分成功和失败。
\end{itemize}

因此，向左失败并不是因为训练数据中多数向左动作都失败。世界模型学到了“成功移动后画面如何变化”，却没有学好“何时根本不会移动”。当外部数据中障碍更多时，这个问题被放大。向前和向右也存在相似但程度不同的错误。

\section{碰撞信息既在状态压缩中损失，也在世界模型中被进一步忽略}

ID71用相同的小型分类器分别读取真实16位置状态、DINO和ID75世界模型预测，判断移动是否成功。

\begin{table}[htbp]
\centering
\caption{ID71动作成功判断AUC}
\small
\begin{tabular}{lccc}
\toprule
\textbf{动作} & \textbf{16位置状态} & \textbf{DINO} & \textbf{世界模型预测} \\
\midrule
向前 & $0.873$ & $0.874$ & $0.738$ \\
向右 & $0.617$ & $0.718$ & $0.540$ \\
向左 & $0.674$ & $0.762$ & $0.617$ \\
\bottomrule
\end{tabular}
\end{table}

这组结果把问题分成了两层：

\begin{enumerate}
  \item 真实状态仍然包含部分动作结果信息，但左右方向弱于DINO；向左状态减去DINO的AUC区间为$[-0.167,-0.013]$，可以确认状态压缩损失了信息。
  \item 世界模型预测又低于直接读取真实状态的分类器，说明世界模型没有充分利用状态中剩余的信息。
\end{enumerate}

因此，仅扩大世界模型不能修复全部问题；仅重做状态也不能自动保证世界模型正确处理动作失败。两层都需要门槛。

\section{从最后隐藏表示修补状态的路线没有成功}

ID191检查了状态压缩之前、与动作生成来自同一次模型计算的隐藏表示。首先，这些隐藏表示经过冻结的SFT1压缩模块能够逐元素重现旧状态，RMSE和最大绝对误差均为0，证明数据链路一致。

随后只训练一个受限的小型残差适配器。它略微改善了与DINO的视觉距离，但没有恢复目标或左右动作结果：

\begin{itemize}
  \item 目标micro准确率：旧状态$6.94\%$，压缩前隐藏表示$5.49\%$，适配后$7.80\%$，DINO$10.98\%$；
  \item 向左AUC：旧状态$0.695$，压缩前隐藏表示$0.625$，适配后$0.622$，DINO$0.778$；
  \item 向右AUC：适配后$0.726$，没有形成对旧状态的稳定改善。
\end{itemize}

这项调查否定的是“只读取最后隐藏表示，再加一个低容量适配器”的具体路线。它不能证明隐藏表示中绝对不存在任何复杂的非线性信息，但继续增加适配器维度、训练轮数或学习率缺乏证据支持。

\section{ID192：任务目标和视觉信息的位置}

ID192在同一次冻结主模型计算中读取四种特征：

\begin{enumerate}
  \item 语言模型之前的当前图片表示；
  \item 经过语言模型处理后的当前图片表示；
  \item instruction词元的输入表示；
  \item instruction词元经过语言模型后的表示。
\end{enumerate}

它同时验证最后隐藏表示能够严格重现ID191状态，RMSE和最大绝对误差均为0。

\subsection{任务目标在instruction中几乎完整存在}

\begin{table}[htbp]
\centering
\caption{ID192目标分类结果}
\small
\begin{tabular}{lcc}
\toprule
\textbf{输入表示} & \textbf{micro top-1} & \textbf{macro top-1} \\
\midrule
instruction输入表示 & \good{$99.71\%$} & \good{$99.00\%$} \\
instruction最终表示 & \good{$99.13\%$} & \good{$98.27\%$} \\
经过语言模型的当前图片表示 & $6.94\%$ & $3.59\%$ \\
最后16位置隐藏表示 & $5.49\%$ & $3.58\%$ \\
现有状态 & $6.07\%$ & $4.26\%$ \\
DINO & $5.78\%$ & $4.04\%$ \\
\bottomrule
\end{tabular}
\end{table}

目标信息在原始instruction表示中非常清楚，但到最后16位置表示和状态时几乎消失。这是本轮调查中最明确的定位结果。新的统一状态应当直接保留instruction表示，而不应期待旧16位置隐藏表示自己恢复目标。

\subsection{最有希望的视觉来源是语言模型处理后的当前图片表示}

\begin{table}[htbp]
\centering
\caption{ID192动作成功判断AUC}
\small
\begin{tabular}{lcccc}
\toprule
\textbf{动作} & \textbf{语言模型前图片} & \textbf{语言模型后图片} & \textbf{最后隐藏表示} & \textbf{DINO} \\
\midrule
向前 & $0.835$ & $0.868$ & $0.865$ & $0.873$ \\
向右 & $0.525$ & $0.740$ & $0.764$ & $0.719$ \\
向左 & $0.718$ & $0.731$ & $0.599$ & $0.738$ \\
\bottomrule
\end{tabular}
\end{table}

语言模型之前的图片表示在向右动作上接近随机。语言模型处理后的当前图片表示在三个方向上都接近DINO，并在向右的观测值上略高于DINO。这表明语言模型对图片和历史上下文的处理很重要，该表示是目前最值得继续验证的视觉来源。

严格门槛仍未通过。外部验证中的向右和向左样本只有142和193条。语言模型后图片表示减去DINO AUC的95\%区间分别为：

\begin{equation}
  [-0.081,+0.125],\qquad [-0.081,+0.067].
\end{equation}

区间包含了“几乎相同”和“有一定差距”两种可能。因此当前证据支持它作为首选候选，但尚不足以证明它“不比DINO差超过0.02”。不能因为这个保守门槛没有通过，就直接断言必须部署第二个DINO视觉模型。

\section{综合诊断}

近期证据可以用三层问题概括。

\subsection{第一层：目标在状态压缩时丢失}

instruction表示的目标分类接近$100\%$，最后隐藏表示和现有状态只有约$5\%$--$7\%$。目标并非没有出现在输入中，而是在生成16位置表示的过程中没有被保留下来。

\subsection{第二层：左右碰撞信息被削弱}

DINO和语言模型处理后的当前图片表示都能较好地区分向左、向右移动是否成功。现有状态和最后隐藏表示在部分方向上明显较弱。视觉信息并非全部消失，但压缩过程没有稳定保留规划所需的侧向障碍信息。

\subsection{第三层：世界模型没有正确处理动作失败}

即使真实状态仍有部分动作结果信息，ID75预测也低于直接读取真实状态的小分类器。世界模型尤其容易把被挡住的动作预测为成功移动，并受训练与外部验证失败率变化影响。

\subsection{目前没有证据支持的解释}

\begin{itemize}
  \item 没有证据表明部署主模型整体损坏；
  \item 没有证据表明只要扩大旧适配器就能恢复信息；
  \item 没有证据表明DINO可以直接替代统一状态，因为DINO不知道instruction；
  \item 没有证据支持继续复用旧ID75世界模型或旧ValueHead；状态定义变化后，它们面对的是不同输入；
  \item 没有足够证据把视觉和目标拆成两个独立部署状态。项目仍应输出一个统一的16位置状态。
\end{itemize}

\section{当前SFT1监督与建议监督的对比}

\subsection{当前SFT1实际在学习什么}

历史SFT1的16位置状态训练代码使用成功轨迹中的单步前缀。16个状态查询词元本身不计算词元预测损失；它们通过后续回答/动作损失和DINO视觉损失得到梯度。训练目标可以概括为：

\begin{equation}
  L_{\mathrm{current}}
  =L_{\mathrm{response/action}}^{\mathrm{successful\ demonstrations}}
  +\lambda_d\frac{1}{16}\sum_{i=1}^{16}
  \left\|z_{t,i}-d_{t,i}^{\mathrm{DINO}}\right\|_2^2.
\end{equation}

视觉模块保持冻结，语言模型使用LoRA这种低参数量更新方法，16个查询词元有可训练的附加向量，\code{SharedSlotProjector}把每个查询位置映射到1024维。DINO目标来自当前图片的$4\times4$区域表示，状态被逐位置、逐元素直接拉向DINO。

这个目标适合得到视觉网格，却没有明确要求状态保存instruction，也没有要求状态判断移动会不会被挡住。只使用成功示范做回答/动作模仿，也不能充分覆盖失败动作。更重要的是，直接惩罚$z-DINO$会使任何DINO没有的额外信息都增加损失，不适合作为完整视觉--目标状态的唯一表示目标。

\subsection{逐项对比}

\begin{longtable}{p{3.0cm}p{5.6cm}p{6.0cm}}
\toprule
\textbf{方面} & \textbf{当前SFT1} & \textbf{建议的新SFT1} \\
\midrule
\endfirsthead
\toprule
\textbf{方面} & \textbf{当前SFT1} & \textbf{建议的新SFT1} \\
\midrule
\endhead
状态是否继承旧状态 & 训练新查询表示和压缩模块，但最终状态被直接拟合到DINO & 不把旧SFT1或ID74状态作为目标；旧状态只作比较基准 \\
视觉监督 & $z$逐元素拟合DINO网格 & 从完整状态经共享的小型线性读取模块恢复DINO方向和16个区域间的空间关系 \\
任务指令 & 没有独立监督 & 匹配冻结ID176从exact instruction得到的固定表示；目标类别分类只作检查或小型辅助 \\
动作被挡 & 没有独立监督 & 从$(z,a)$预测动作是否成功；优先为同一模拟器状态采集所有主要移动动作的结果 \\
目标与画面关系 & 没有可靠监督 & 重新采集目标是否可见；可见时监督左/中/右位置和大致距离 \\
训练数据 & 回答/动作模仿主要使用成功轨迹 & 视觉和instruction使用所有pre-RL观察；动作结果使用全部转移并保留失败；只有模仿损失使用成功示范 \\
策略保护 & 回答/动作损失同时改变状态和策略 & 用冻结ID176的动作分布约束新查询词元，不允许状态改善以actor退化为代价 \\
checkpoint选择 & 主要看总训练/验证损失 & 视觉、目标、逐动作成功判断和策略保持分别过门槛；不允许一项高分掩盖另一项失败 \\
\bottomrule
\end{longtable}

\subsection{建议的新监督信号}

训练时允许增加几个容量很小的读取模块。它们从同一个完整状态读取信息，只用于监督和检查，部署时删除；这不会把状态拆成视觉和目标两个分支。

\paragraph{视觉空间。}
令$R_v$为16个位置共用的线性读取模块，$d_{t,i}$为原始观察图片的DINO区域表示：

\begin{equation}
  \hat d_{t,i}=R_v(z_{t,i}),
\end{equation}

定义$S(x)_{ij}=\cos(x_i,x_j)$，即任意两个区域之间的相似度。视觉损失为：

\begin{equation}
  L_{\mathrm{visual}}
  =\frac{1}{16}\sum_i\left[1-\cos(\hat d_{t,i},d_{t,i})\right]
  +\eta\left\|S(\hat d_t)-S(d_t)\right\|_F^2.
\end{equation}

第一项保存每个区域的视觉内容，第二项保存16个区域之间的空间关系。它们不要求整个状态等于DINO。

\paragraph{任务指令。}
用冻结ID176对完整、真实instruction生成固定目标：

\begin{equation}
  g_t^{\mathrm{teacher}}
  =\operatorname{stopgrad}\!\left(E_{\mathrm{ID176}}(\mathrm{instruction}_t)\right),
\end{equation}

\begin{equation}
  L_{\mathrm{instruction}}
  =1-\cos\!\left(
  R_g\!\left(\frac{1}{16}\sum_i z_{t,i}\right),
  g_t^{\mathrm{teacher}}
  \right).
\end{equation}

教师表示必须预先计算并固定，不能与新模型一起移动。ID192已经验证该表示能够可靠读出目标。目标类别准确率继续作为外部检查，避免只匹配无关措辞。

\paragraph{动作可行性。}
理想数据应在同一个模拟器状态的独立副本中分别尝试向前、向左和向右，得到每个动作的真实成功标签$y_t^{(a)}$：

\begin{equation}
  L_{\mathrm{feasibility}}
  =\sum_{a\in\mathcal A_{\mathrm{move}}}
  \operatorname{BCE}\!\left(R_a(z_t,a),y_t^{(a)}\right).
\end{equation}

现有归档只有实际执行动作的下一条环境反馈，可用于较弱的小规模验证，不能替代全动作标签。真实结果只能作为训练目标，绝不能作为部署输入。

\paragraph{目标与画面的关系。}
正式状态还应监督目标当前是否可见；仅在可见时监督目标在画面的左/中/右位置和大致距离。看不见时不强制预测全局最短距离，因为该信息可能无法从现有观察和历史推断。现有归档的逐行任务身份不可靠，这部分需要新的、身份严格绑定的pre-RL采集。

\paragraph{策略保持。}
改变状态查询词元可能改变后续动作概率。应固定ID176教师动作分布，并约束：

\begin{equation}
  L_{\mathrm{policy}}
  =D_{\mathrm{KL}}\!\left(
  \pi_{\mathrm{ID176}}(a\mid o_t,\mathrm{instruction},\mathrm{history})
  \;\|\;
  \pi_{\mathrm{new}}(a\mid o_t,\mathrm{instruction},\mathrm{history})
  \right).
\end{equation}

如果后续解冻Qwen参数，还要保留原始成功示范的回答/动作损失，防止CoT、格式和动作能力退化。

\subsection{建议目标及其边界}

正式状态训练目标为：

\begin{equation}
  L_{\mathrm{state}}
  =\lambda_vL_{\mathrm{visual}}
  +\lambda_gL_{\mathrm{instruction}}
  +\lambda_aL_{\mathrm{feasibility}}
  +\lambda_rL_{\mathrm{relation}},
\end{equation}

\begin{equation}
  L_{\mathrm{total}}
  =L_{\mathrm{state}}+\lambda_pL_{\mathrm{policy}},
\end{equation}

若解冻Qwen，再加入$\lambda_{\mathrm{SFT}}L_{\mathrm{SFT}}$。各损失先按样本平均并检查其对状态查询位置产生的梯度大小，避免某一项比其他项大几十倍。checkpoint仍按各项外部门槛分别选择，不按加权总分选择。

不应作为新状态监督的信号包括：旧SFT1或ID74状态、旧世界模型预测、旧ValueHead/Q、trajectory reward、旧decoder重建、图片变化代理、fixed CoT及任何post-RL heldout数据。

\subsection{当前可用数据与缺口}

\begin{table}[htbp]
\centering
\caption{建议监督信号的数据准备情况}
\small
\begin{tabularx}{\textwidth}{p{4.0cm}p{2.4cm}X}
\toprule
\textbf{信号} & \textbf{当前是否可用} & \textbf{说明} \\
\midrule
原始图片DINO网格 & 是 & 已有冻结cache，但正式新状态应重新锁定exact source和原图处理身份 \\
冻结instruction表示 & 是 & ID192已有3,566条record级固定表示，可覆盖现有train/validation \\
实际执行动作的成功/失败 & 是 & 来自下一条精确环境反馈；有策略选择偏差，只适合小规模验证 \\
同一状态的全动作可行性 & 否 & 需要在模拟器状态副本中逐动作测试并保存 \\
目标可见性及可见时位置 & 否 & 需要逐行任务身份正确的新pre-RL采集 \\
冻结ID176动作分布 & 可只读生成 & 需要从同一次前缀计算并固定，不能使用更新后的模型重算教师 \\
\bottomrule
\end{tabularx}
\end{table}

\section{建议的后续路线}

\subsection{第一步：扩大只读统计检查}

ID192已经保存14,261个状态的冻结特征缓存。先进行按原始图片内容分组的五折交叉验证，缩小“语言模型后图片表示与DINO相差多少”的统计范围。该检查不能冒充新的外部验证；现有归档外部验证仍是主要外部结果。

\subsection{第二步：使用现有数据做状态监督小规模验证}

冻结Qwen和视觉模块，只训练16个查询词元的附加参数、新\code{SharedSlotProjector}和训练用读取模块。使用：

\begin{equation}
  L_{\mathrm{visual}}+L_{\mathrm{instruction}}
  +L_{\mathrm{actual\ action\ result}}+L_{\mathrm{policy}}.
\end{equation}

这一步只验证查询词元能否主动读取ID192定位的图片和instruction信息。由于缺少全动作与可靠目标关系标签，它不能直接建立正式统一状态。

\subsection{第三步：重新采集正式监督数据}

在身份严格绑定的pre-RL模拟器状态中保存所有主要移动动作的可行性、目标是否可见，以及可见时的位置和距离。失败与被挡转移必须保留。视觉和instruction监督可使用所有观察；只有回答/动作模仿损失使用成功示范。

\subsection{第四步：正式重训SFT1状态提取}

使用完整四类状态监督和策略保持约束。先保持Qwen冻结；如果查询词元仍无法读取信息，再只解冻语言模型上层少量参数并加入原始SFT损失。现有SFT1状态只作基准，不作为新状态输入或训练目标。

\subsection{第五步：从头训练能表示动作失败的一步世界模型}

新世界模型应预测动作成功概率，并分别建模“成功移动后的状态”和“动作被挡住后的状态”。一步模型必须按每个主要动作，以及成功和被挡住子集分别优于复制基准。通过后才扩展到第2步和第4步。

\subsection{第六步：重新训练价值模型并恢复搜索}

新状态定义确定后，从头训练并校准ValueHead和动作价值。先检查真实状态和一步预测状态，再检查第2--4层预测。只有预测深度增加时价值排序和回报校准没有持续恶化，才恢复K4搜索和新的联合强化学习。

\section{尚未解决的问题}

\begin{enumerate}
  \item 语言模型处理后的当前图片表示能否在更大、按图片分组的数据上稳定达到DINO水平？
  \item 目标与当前画面的关系---例如目标是否可见、位于哪个方向、任务进行到哪一步---能否由新的统一状态读出？接近$100\%$的目标类别识别只证明instruction内容存在，不等于完整任务理解。
  \item ValueHead在第2--4层预测状态上的均值、方差、动作排序和回报校准如何变化？
  \item 当前归档缺少足够多的“相同图片、不同真实目标”样本，因此还不能做强目标因果检查。
  \item ID189终止记录没有精确的最终16位置状态，相关诊断只能覆盖非终止转移。
\end{enumerate}

\section{可复现记录}

\begin{longtable}{p{1.3cm}p{2.2cm}p{2.3cm}p{8.3cm}}
\toprule
\textbf{调查} & \textbf{Slurm Job} & \textbf{结果} & \textbf{关键产物或说明} \\
\midrule
\endfirsthead
\toprule
\textbf{调查} & \textbf{Slurm Job} & \textbf{结果} & \textbf{关键产物或说明} \\
\midrule
\endhead
ID58 & 528812 & 完成 & checkpoint状态矩阵；定位压缩模块漂移。 \\
ID59 & 528906 & 完成 & 部署主模型与SFT1视觉检查。 \\
ID60 & 528931 & 完成 & 冻结目标分类；cache SHA256：\code{0fa99413...f0daf8b6}。 \\
ID75 & 529411 & 完成，门槛失败 & 一步模型总体改善，但向左失败；result SHA256：\code{e56bfdcd...16958a796}。 \\
ID61 & 529546 & 完成 & 动作成功/被挡分层；result SHA256：\code{bace6fcb...105a28e4}。 \\
ID71 & 529619 & 完成 & 状态动作结果分类；result SHA256：\code{2e2e1675...b70c113}。 \\
ID191 & 529703 & 完成，门槛失败 & 最后隐藏表示小适配器；result SHA256：\code{1e1307c2...506bb}。 \\
ID192 & 529879 & 完成，部分门槛通过 & 目标来源明确；视觉候选仍受统计范围限制。feature cache SHA256：\code{a46278d4...f813c}；result SHA256：\code{4592b68d...e7717a}。 \\
\bottomrule
\end{longtable}

ID192正式尝试0因词元边界处理错误而在特征生成前失败；retry1完成所有模型计算后因批次数组拼接过慢被人工取消。两次都没有生成可用分类器、checkpoint或参数更新。retry2改为按最终位置预分配连续数组后成功完成。这些失败不改变科学结论，但已记录以避免重复。

\section{结论}

近期调查已经把问题从笼统的“长期规划不好”缩小到可以行动的状态接口问题：

\begin{enumerate}
  \item 部署主模型基本健康，旧状态压缩模块曾发生严重漂移；
  \item 即使使用较健康的SFT1压缩模块，当前16位置状态仍几乎丢失instruction目标，并削弱部分左右碰撞信息；
  \item 旧世界模型进一步没有充分使用剩余信息，尤其不能可靠预测动作被挡住；
  \item 目标信息已经明确定位在instruction表示中；
  \item 经过语言模型处理后的当前图片表示是最有希望的视觉来源，但还需要更大分组检查来缩小不确定范围；
  \item 当前SFT1的“成功示范回答/动作损失加状态直接拟合DINO”没有覆盖任务指令、动作被挡和目标--画面关系，监督目标需要重做。
\end{enumerate}

在这些证据下，继续增加旧世界模型训练轮数、扩大ID191适配器，或直接恢复K4强化学习，都不能解决基础输入问题。合理顺序是：扩大冻结统计检查，用现有数据验证新的监督能否驱动16个查询位置读取图片和instruction；随后补采全动作可行性与目标关系标签，正式重训SFT1状态提取；新状态通过所有独立门槛后，再从头训练一步世界模型、长预测、价值模型和搜索策略。

\appendix
\section{内部资料索引}

本报告数字来自以下项目文件和只读产物：

\begin{itemize}
  \item \url{ai_tasks/wm_visual_goal_state_alignment_plan.md}
  \item \url{ai_tasks/ai_progress/2026-08-23_frozen_sft1_goal_probe_residual_t1_canary.md}
  \item \url{ai_tasks/id191_state_interface_direction_canary_contract.md}
  \item 历史SFT1 DINO-grid监督实现：commit \code{6c1828a09615d8fc9eb9ac77b5723eefa1c6eb2e}中的\code{experiments/training/sft1/train\_dino\_grid.py}与\code{src/nimloth/training/sft1/dino\_grid.py}
  \item \url{ai_tasks/id192_multimodal_feature_location_audit_contract.md}
  \item \url{/workspace/remote2/nimloth-artifacts/id61_id75_action_outcome_audit_retry1/result.json}
  \item \url{/workspace/remote2/nimloth-artifacts/id71_frozen_state_action_outcome_probe/result.json}
  \item \url{/workspace/remote2/nimloth-artifacts/id191_state_interface_direction_canary_retry1/result.json}
  \item \url{/workspace/remote2/nimloth-artifacts/id192_multimodal_feature_location_audit_retry2/result.json}
\end{itemize}

这些资料同时保留了完整配置、split限制、置信区间、产物哈希和失败尝试记录。本报告为人类阅读而压缩了实现细节；若表格中的四舍五入数字与JSON存在末位差异，以对应\code{result.json}为准。

\end{document}
